\documentclass[11pt]{article}
\usepackage[margin=1in]{geometry}
\usepackage[T1]{fontenc}
\usepackage{lmodern}
\usepackage{amsmath,amssymb,amsthm}
\usepackage{graphicx}
\usepackage[hidelinks]{hyperref}
\usepackage{microtype}
\usepackage{enumitem}

\newtheorem*{theorem}{Theorem}
\newtheorem*{lemma}{Bridge lemma}
\newcommand{\N}{\mathbb N}
\newcommand{\K}{\mathbb K}
\newcommand{\one}{\mathbf 1}

\begin{document}

\begin{center}
{\Large Full solution packet for arXiv:2205.06839}\\[4pt]
{\large Arbitrary-weight disjoint superdemocracy implies superdemocracy}
\end{center}

\paragraph{Status.}
\textbf{full\_solution\_likely\_valid}. This packet answers Question 2 in
Section 6 of H\`ung Vi\d{\^e}t Chu, \emph{On weighted greedy-type bases},
arXiv:2205.06839 \cite{Chu}.

\paragraph{Source question.}
In Definition 1.2, a set weight is an arbitrary map
\(w:\mathcal P(\N)\to[0,\infty]\) such that \(w(\varnothing)=0\) and
\(w(A)>0\) for every nonempty \(A\). Question 2 on source page 14 asks:

\begin{quote}\small
For weights in Definition 1.2, is an \(\omega\)-disjoint superdemocratic
basis also \(\omega\)-superdemocratic? If not, what minimal condition(s) to
put on \(\omega\) so that the two properties are equivalent?
\end{quote}

\begin{center}
\includegraphics[width=.98\linewidth]{figures/open_problem_crop.png}
\end{center}

\paragraph{Definitions.}
Let \(\mathcal B=(e_n)_{n=1}^{\infty}\) be the semi-normalized Schauder basis
fixed in the source paper, with biorthogonal functionals \((e_n^*)\). Put
\[
 M=\sup_n\|e_n\|<\infty,
 \qquad
 M^*=\sup_n\|e_n^*\|<\infty.
\]
For a finite set \(A\) and signs \(\varepsilon=(\varepsilon_n)_{n\in A}\),
where \(|\varepsilon_n|=1\), write
\[
 \one_{\varepsilon A}=\sum_{n\in A}\varepsilon_n e_n.
\]
The basis is \(K\)-\(\omega\)-disjoint superdemocratic if
\[
 \|\one_{\varepsilon A}\|
 \le K\|\one_{\delta B}\|
 \tag{1}\label{eq:disjoint}
\]
whenever \(A,B\subset\N\) are finite, \(A\cap B=\varnothing\),
\(w(A)\le w(B)\), and the signs are arbitrary. It is
\(\omega\)-superdemocratic when the same comparison holds without the
disjointness requirement.

\paragraph{Proof intuition.}
The weight itself need not supply a fresh set of intermediate weight, so the
source proof for structured weights cannot be copied. Instead, we make a
fresh set of intermediate \emph{norm}. If signed constant-coefficient sums
are unbounded, then they remain unbounded after deleting any finite set of
coordinates. Adding their terms one at a time gives a signed sum whose norm
lies just above any prescribed level. For two putative counterexample sums
of norms \(a\gg b\), choose a fresh sum of norm strictly between \(Kb\) and
\(a/K\). The disjoint comparison with the first set forces the fresh set's
weight below the first weight. Transitivity places it below the second
weight, and the disjoint comparison with the second set then contradicts
its norm being greater than \(Kb\). If all signed sums are bounded instead,
the conclusion follows immediately from the uniform lower bound on every
nonzero signed sum.

\begin{lemma}
Suppose
\[
 \sup\{\|\one_{\varepsilon E}\|:
 E\subset\N\text{ finite},\ |\varepsilon_n|=1\}=\infty.
 \tag{2}\label{eq:unbounded}
\]
For every finite \(F\subset\N\) and every \(t\ge0\), there are a nonempty
finite \(C\subset\N\setminus F\) and signs \(\eta\) such that
\[
 t<\|\one_{\eta C}\|\le t+M.
 \tag{3}\label{eq:bridge}
\]
\end{lemma}

\begin{proof}
First, signed sums supported in \(\N\setminus F\) are unbounded. Indeed, if
they were bounded by \(L\), then every signed sum would satisfy
\[
 \|\one_{\varepsilon E}\|
 \le \sum_{n\in E\cap F}\|e_n\|
      +\|\one_{\varepsilon(E\setminus F)}\|
 \le |F|M+L,
\]
contrary to \eqref{eq:unbounded}.

Choose a finite \(D\subset\N\setminus F\) and signs \(\eta\) such that
\(\|\one_{\eta D}\|>t+M\). Enumerate
\(D=\{n_1,\ldots,n_m\}\) and let
\(z_j=\sum_{i=1}^j\eta_{n_i}e_{n_i}\), with \(z_0=0\). Let \(j\) be the
first index for which \(\|z_j\|>t\). Then \(\|z_{j-1}\|\le t\), and hence
\[
 t<\|z_j\|
 \le \|z_{j-1}\|+\|e_{n_j}\|
 \le t+M.
\]
Taking \(C=\{n_1,\ldots,n_j\}\) proves the claim.
\end{proof}

\begin{theorem}[Affirmative answer to Question 2]
Let \(w:\mathcal P(\N)\to[0,\infty]\) be any weight in Definition 1.2 of
the source paper. Every \(\omega\)-disjoint superdemocratic Schauder basis is
\(\omega\)-superdemocratic. No additional condition on \(w\) is required.

More precisely, suppose \eqref{eq:disjoint} holds with constant \(K\), and
write
\[
 D=\sup\{\|\one_{\varepsilon E}\|:
 E\subset\N\text{ finite},\ |\varepsilon_n|=1\}.
\]
If \(D<\infty\), a superdemocracy constant is \(DM^*\). If \(D=\infty\),
a superdemocracy constant is
\[
 K^2+KMM^*.
 \tag{4}\label{eq:constant}
\]
\end{theorem}

\begin{proof}
Let \(A,B\subset\N\) be finite, let \(w(A)\le w(B)\), and fix signs
\(\varepsilon,\delta\). We must bound
\[
 a=\|\one_{\varepsilon A}\|
 \quad\text{by a constant times}\quad
 b=\|\one_{\delta B}\|.
\]
If \(A=\varnothing\), there is nothing to prove. If \(A\ne\varnothing\),
then \(B\ne\varnothing\), because \(w(A)>0=w(\varnothing)\). For any
\(j\in B\), biorthogonality gives
\[
 1=|e_j^*(\one_{\delta B})|
 \le \|e_j^*\|b\le M^*b,
 \qquad\text{so}\qquad b\ge (M^*)^{-1}.
 \tag{5}\label{eq:lower}
\]

If \(D<\infty\), then \(a\le D\le DM^*b\), proving the claim in this
case. Notice that this argument does not use the weight at all.

Now assume \(D=\infty\). We claim that
\[
 a\le (K^2+KMM^*)b.
\]
Suppose otherwise. By \eqref{eq:lower},
\[
 a>(K^2+KMM^*)b
   \ge K^2b+KM
   =K(Kb+M).
 \tag{6}\label{eq:room}
\]
Apply the bridge lemma to \(F=A\cup B\) and \(t=Kb\). We obtain a finite
\(C\subset\N\setminus(A\cup B)\) and signs \(\eta\) such that, on writing
\(c=\|\one_{\eta C}\|\),
\[
 Kb<c\le Kb+M<a/K,
 \tag{7}\label{eq:middle}
\]
where the last inequality is \eqref{eq:room}.

If \(w(A)\le w(C)\), then \(A\cap C=\varnothing\) and
\eqref{eq:disjoint} would give \(a\le Kc<a\), a contradiction. Since the
extended real numbers are totally ordered, we must therefore have
\[
 w(C)<w(A)\le w(B).
\]
But \(C\cap B=\varnothing\), so another application of
\eqref{eq:disjoint} gives \(c\le Kb\), contradicting
\eqref{eq:middle}. This proves \eqref{eq:constant} and completes the proof.
\end{proof}

\paragraph{Consequences.}
The answer to the first sentence of Question 2 is yes. Therefore the second
sentence has the strongest possible answer: within Definition 1.2, there are
no additional conditions to impose on \(\omega\). Proposition 6.1 of the
source paper, which proves the implication for structured weights, follows
as a special case.

The proof in fact needs less than the source's Schauder-basis hypotheses. It
uses neither convergence of basis expansions nor boundedness of partial-sum
projections. Consequently, the same result holds for every biorthogonal system
\((e_n,e_n^*)\) for which both \(\sup_n\|e_n\|\) and
\(\sup_n\|e_n^*\|\) are finite.

\paragraph{Verification notes.}
The proof uses only four structural facts: finite supports have an infinite
complement; \(\sup_n\|e_n\|<\infty\); \(\sup_n\|e_n^*\|<\infty\); and the
total order on \([0,\infty]\). Thus weights equal to \(\infty\) cause no
exception. The bounded case and unbounded case exhaust all possibilities.
The bridge construction is valid over both \(\mathbb R\) and \(\mathbb C\),
since multiplying a basis vector by a scalar of modulus one does not alter
its norm.

The run's lightweight indexes and local arXiv source corpus were searched
for arXiv:2205.06839, the title, the exact Question 2 wording, and close
variants of ``disjoint superdemocratic'' with arbitrary set weights. Bounded
web and forward-citation searches found no exact later solution. Crossref and
OpenAlex returned no citing works for the published DOI; Semantic Scholar's
sole result was checked directly and contains only a commented-out reference,
with no discussion of Question 2. This is a bounded rather than exhaustive
search, so novelty confidence is moderate-to-high pending expert review; full
details are in \texttt{citation\_audit.md}.

\paragraph{Limitations.}
This packet answers Question 2 only. It does not address Question 1 about
extending the equivalence of \(\omega\)-almost greedy and
\(\omega\)-semi-greedy bases beyond structured weights. The quantitative
constant in the bounded-sum case depends on \(D\); no optimality claim is
made for either displayed constant.

\paragraph{Human-review recommendation.}
Review as a candidate full solution. The most important checks are (i) the
claim that unbounded signed sums remain unbounded after deleting finitely many
coordinates, and (ii) the two weight-order deductions following
\eqref{eq:middle}. Both are explicit in the proof.

\begin{thebibliography}{9}
\bibitem{Chu}
H. V. Chu,
\emph{On weighted greedy-type bases},
Bulletin of the Brazilian Mathematical Society, New Series \textbf{54}
(2023), Article 49; arXiv:2205.06839; DOI:10.1007/s00574-023-00367-3.
\end{thebibliography}

\end{document}
